{\footnotesize
\begin{xltabular}{\linewidth}{l p{0.28\linewidth} r X}
\caption{Registro de decisiones de depuración}\label{tab:decisiones} \\
\toprule
Id. & Decisión & Casos & Acción \\
\midrule
\endfirsthead
\multicolumn{4}{l}{\emph{(continuación)}} \\
\toprule
Id. & Decisión & Casos & Acción \\
\midrule
\endhead
\midrule
\multicolumn{4}{r}{\emph{(continúa)}} \\
\endfoot
\bottomrule

\endlastfoot
D01 & Geography separada en condado y estado & \num{3047} & Se separa por la última coma en Geography (condado) y state (estado). \\
D02 & Notificadomuerte excluida por fuga de información & \num{2841} & Se conserva en los datos pero se excluye de todos los modelos. \\
D03 & Valor centinela en la incidencia & \num{206} & Se sustituye por valor ausente. \\
D04 & Edad mediana registrada en meses & \num{30} & Se divide entre 12 y se verifica que el resultado quede entre las medianas masculina y femenina. \\
D05 & Tamaño medio del hogar dividido por 100 & \num{61} & Se multiplica por 100. \\
D06 & PctSomeCol18\_24 recuperada por identidad contable & \num{2285} & Se calcula como 100 \textminus{} (sin bachillerato + bachillerato + grado). \\
D07 & Población nativa o multirracial (variable derivada) & \num{86} & Se crea PctNativeMulti = 100 \textminus{} (blanca + negra + asiática + otra). \\
D08 & Incidencia inverosímil (mortalidad $\geq$ incidencia) & \num{1} & Se marca la incidencia como ausente. \\
D09 & Región censal & \num{3047} & Se agrupan en las cuatro regiones oficiales de la Oficina del Censo. \\
D10 & Transformaciones de variables muy asimétricas & \num{3047} & Población: logaritmo; renta: logaritmo; ensayos: log1p. \\
D11 & Variables excluidas de los modelos & \num{12} & Se excluyen: TARGET\_deathRate, avgDeathsPerYear, avgAnnCount, Notificadomuerte, binnedInc, Geography, state, MedianAgeMale, MedianAgeFemale, PctPrivateCoverageAlone, PctSomeCol18\_24, PctWhite. \\
\end{xltabular}
}
